\section{Conclusion}\label{sec:conclusion}

We have completed the complexity map of FOND HTN planning for the solution concepts of
\CBa{} and \CBb{} and the standard structural classes. The new entries confirm the general
pattern that outcome dependence costs one exponential, but with notable exceptions. Strong
outcome-dependent solutions with fixed methods become undecidable as soon as the hierarchy allows
two independent recursive chains. They are therefore strictly harder than method-based policies
on regular and tail-recursive problems, because the planner must commit to the refinement before
it observes anything. For strong cyclic method-based policies the right definition requires
fairness, and with it the unrestricted problem leaves the arithmetical hierarchy altogether.

Several questions remain. The complexity of strong cyclic solutions under the literal reading of
\CBb{} is still open. It is a well-defined, if unintended, notion, and its winning condition is a
pure safety condition, so we expect it to be easier than the fair reading. Our \EXPTIME{} and
\TWOEXP{} lower bounds for method-based policies use negative preconditions. They can presumably be
compiled away as in classical planning, but this has not been checked formally. Finally, a formal
proof of Harel's theorem would make the last entry of Table~\ref{tab:results} fully
machine-checked as well.
